\documentclass[a4paper,12pt]{article}
\usepackage[utf8]{inputenc}
\usepackage[italian]{babel}
\usepackage[T1]{fontenc}
\usepackage{enumitem}
\usepackage{geometry}
\usepackage{listings}
\usepackage{xcolor}
\usepackage{tcolorbox} % Aggiunto per i riquadri
\usepackage{hyperref}  % Aggiunto per i link

\geometry{margin=2.5cm}

\title{Tutorato Programmazione 2}
\author{Michele Porcellato, Giacomo Vettore}
\date{15 Aprile 2026}

\begin{document}

\maketitle

\section{Note}
In caso di dubbi, consultare la \href{https://docs.oracle.com/en/java/javase/25/docs/api/index.html}{\textit{documentazione Java}}. Sarà accessibile anche durante l'esame.

\section{Esercizi}

\subsection{QuestLog}
Si progetti \textbf{QuestLog}, un'applicazione dove selezionare e risolvere diverse Quests per Geraldo da Ravina (TN).

Tutte le Quests hanno un nome, un livello minimo, una tipologia e una ricompensa, cioè la quantità di denaro che si guadagna risolvendole. Esistono le seguenti tipologie di Quests, che aggiungono caratteristiche specifiche:

\begin{itemize}
    \item \textbf{Main}: indicano anche il Luogo dove avvengono. Tutte le Main Quests danno 10 di ricompensa. Inoltre, ognuna di queste Quest indica la sua prossima, a meno che non sia la Quest finale (la prossima è \texttt{null}). Per risolvere una Main Quest bisogna aver risolto la sua precedente.
    \item \textbf{Secondary}: indicano un non-playing character (NPC) con cui interagire.
    Risolvere una Secondary Quest aumenta il livello di 2 invece che di 1. 
    \item \textbf{Contract}: indicano la lista di mostri da sconfiggere. Quando sono risolte, le Contract Quests offrono il 20\% in più della ricompensa indicata. 
    \item \textbf{Treasure Hunt}: Un sottoinsieme delle Quest Contracts che specificano il Luogo dove trovare del loot dopo aver sconfitto i mostri.
    \item \textbf{Downloadable Content (DLC)}: sono come le Main (definiscono un Luogo e hanno una sequenza, danno ricompensa 10), ma, inoltre, indicano una Main Quest da cui dipendono: se non è stata risolta quella Main Quest, non si può affrontare la specifica DLC Quest. 
\end{itemize}

Inizializzate l'applicazione con le seguenti Quests: 

\begin{itemize}
    \item \textbf{Main Quest}: Lilac and Gooseberries. LV: 2. Location: White Orchard. Prossima: The Nilfgaardian Connection.
    \item \textbf{Main Quest}: The Nilfgaardian Connection. LV: 3. Location: Velen. Prossima: Pyres of Novigrad.
    \item \textbf{Main Quest}: Pyres of Novigrad. LV: 5. Location: Novigrad. Prossima: nessuna.
    \item \textbf{Secondary Quest}: A Favor for a Friend. LV: 2. Ricompensa: 10. NPC: Keira Metz.
    \item \textbf{Secondary Quest}: The Last Wish. LV: 6. Ricompensa: 30. NPC: Bloody Baron.
    \item \textbf{Contract}: Swamp Thing. LV: 12. Ricompensa: 40. Mostri: Drowners, Foglet.
    \item \textbf{Contract-Treasure Hunt}: Coast of Wrecks. LV: 4. Ricompensa: 50. Mostri: Drowners. Luogo: Novigrad.
    \item \textbf{Contract-Treasure Hunt}: Dirty Funds. LV: 1. Ricompensa: 20. Mostri: Wolves, Drowners. Luogo: Velen.
    \item \textbf{DLC Quest}: Envoys, Wineboys. LV: 8. Location: Velen. Prossima: Capture the Castle. Dipende da: The Nilfgaardian Connection. 
    \item \textbf{DLC Quest}: Capture the Castle. LV: 9. Location: Toussaint. Prossima: nessuna. Dipende da: The Nilfgaardian Connection.
\end{itemize}

In alto viene mostrato il livello corrente di Geraldo (inizializzato a 2) e i suoi soldi (inizializzati a 0). Dopo di che verranno stampate a video le Quests disponibili, numerate. Per selezionare una Quest e completarla, basterà digitare il numero corrispondente. Premere Invio senza scrivere nulla per uscire dall'applicazione.

Completare una Quest aggiunge 1 livello a Geraldo (2 nel caso delle Secondary) e la relativa ricompensa ai suoi soldi. Le Quest completate mostrano \texttt{[COMPLETATA]} accanto al nome.

Se la Quest non può essere completata (livello di Geraldo insufficiente, quest precedente non completata per le Main, Main dipendente non completata per le DLC), viene lanciata e gestita un'eccezione opportuna, e viene stampato un messaggio di errore sulla console. Se si prova a completare una Quest già completata, viene lanciata e gestita un'eccezione apposita.

\begin{itemize}
    \item Si usi, quando evidente/utile, una gerarchia di ereditarietà.
    \item Si usino costanti ove ragionevole, e si badi alla pulizia del codice, evitando duplicazioni e codice inutilmente complesso.
    \item Si richiede il class diagram UML. È sufficiente la vista di alto livello (no attributi/metodi) ma devono essere mostrate le associazioni più importanti, oltre a quella di ereditarietà.
\end{itemize}

\subsection*{Esempio 1}
Stato iniziale. L'utente prova a completare la quest n. 6 (Swamp Thing, LV minimo 12), ma Geraldo è di livello 2.

\begin{tcolorbox}[colback=gray!5, colframe=gray!50]
\texttt{Livello: 2 Denaro: 0} \\
\texttt{Quests disponibili:} \\
\texttt{1) Lilac and Gooseberries Ricompensa: 10} \\
\texttt{\ \ \ Tipo: Main Luogo: White Orchard  Prossima: The Nilfgaardian Connection} \\
\texttt{2) The Nilfgaardian Connection Ricompensa: 10} \\
\texttt{\ \ \ Tipo: Main Luogo: Velen Prossima: Pyres of Novigrad} \\
\texttt{3) Pyres of Novigrad Ricompensa: 10} \\
\texttt{\ \ \ Tipo: Main Luogo: Novigrad Prossima: nessuna} \\
\texttt{4) A Favor for a Friend Ricompensa: 10} \\
\texttt{\ \ \ Tipo: Secondary NPC: Keira Metz} \\
\texttt{5) The Last Wish Ricompensa: 30} \\
\texttt{\ \ \ Tipo: Secondary NPC: Bloody Baron} \\
\texttt{6) Swamp Thing Ricompensa: 40} \\
\texttt{\ \ \ Tipo: Contract Mostri: Drowners, Foglet} \\
\texttt{7) Coast of Wrecks Ricompensa: 50} \\
\texttt{\ \ \ Tipo: Treasure Hunt Mostri: Drowners Luogo: Novigrad} \\
\texttt{8) Dirty Funds Ricompensa: 20} \\
\texttt{\ \ \ Tipo: Treasure Hunt Mostri: Wolves, Drowners Luogo: Velen} \\
\texttt{9) Envoys, Wineboys Ricompensa: 10} \\
\texttt{\ \ \ Tipo: DLC Luogo: Velen Prossima: Capture the Castle} \\
\texttt{10) Capture the Castle Ricompensa: 10} \\
\texttt{\ \ \ \ Tipo: DLC  Luogo: Toussaint Prossima: nessuna} \\
\\
\texttt{Seleziona il numero di una Quest da completare.} \\
\texttt{Premi Invio senza scrivere nulla per uscire.} \\
\texttt{>} \textbf{6} \\
\texttt{Errore: livello insufficiente per Swamp Thing (richiesto: 12, attuale 2).}
\end{tcolorbox}

\newpage

\subsection*{Esempio 2}
Geraldo è al livello 5 con 20 denaro. L'utente prova a completare di nuovo la quest 1, già completata.

\begin{tcolorbox}[colback=gray!5, colframe=gray!50]
\texttt{Livello: 5 Denaro: 20} \\
\texttt{Quests disponibili:} \\
\texttt{1) [COMPLETATA] Lilac and Gooseberries Ricompensa: 10} \\
\texttt{\ \ \ Tipo: Main Luogo: White Orchard Prossima: The Nilfgaardian Connection} \\
\texttt{2) The Nilfgaardian Connection Ricompensa: 10} \\
\texttt{\ \ \ Tipo: Main Luogo: Velen Prossima: Pyres of Novigrad} \\
\texttt{3) Pyres of Novigrad Ricompensa: 10} \\
\texttt{\ \ \ Tipo: Main Luogo: Novigrad Prossima: nessuna} \\
\texttt{4) [COMPLETATA] A Favor for a Friend Ricompensa: 10} \\
\texttt{\ \ \ Tipo: Secondary NPC: Keira Metz} \\
\texttt{...} \\
\\
\texttt{Seleziona il numero di una Quest da completare.} \\
\texttt{Premi Invio senza scrivere nulla per uscire.} \\
\texttt{>} \textbf{1} \\
\texttt{Errore: la quest Lilac and Gooseberries e' gia' stata completata.}
\end{tcolorbox}

\section{Suggerimenti}
La soluzione proposta usa MVC. All'esame sarà richiesto ma per ora è possibile svolgere l'esercizio anche senza.
\newline
Si consiglia di implementare una gerarchia di eccezioni \texttt{QuestException} con sottoclassi specifiche per i diversi casi di errore.

\end{document}